\documentclass[12pt]{article}
\usepackage[utf8]{inputenc}
\usepackage[T1]{fontenc}
\usepackage[margin=1in]{geometry}
\usepackage{setspace}
\usepackage{xcolor}
\usepackage[colorlinks=true, urlcolor=blue!60!black, linkcolor=blue!60!black]{hyperref}
\setlength{\parindent}{0em}
\setlength{\parskip}{0.6em}
\renewcommand{\baselinestretch}{1.0}
\begin{document}
\begin{center}
{\LARGE \textit{Saenz v. Roe}, 526 U.S. 489 (1999)\par}
\vspace{0.5em}
{\large Case Brief --- Two-Page Edition\par}
\end{center}

\section{Memory Jogger}
Reviving the Privileges or Immunities Clause, the Court held 7--2 that California could not cap newcomers' first-year welfare benefits at their prior state's level, and that Congress could not authorize it to (\href{https://www.law.cornell.edu/supremecourt/text/98-97}{526 U.S. 489}).

\section{Detailed Case Facts}
To save \$10.9 million a year, California limited families resident for under twelve months to the welfare grant of their prior state; a mother and child from Arizona, for example, would get \$275 rather than \$456. PRWORA, 42 U.S.C. \S~604(c), expressly authorized such caps, and newcomers Brenda Roe and Anna Doe sued for a class when enforcement began in 1997.

\section{Procedural History}
The district court enjoined the cap, the Ninth Circuit affirmed, and the Supreme Court granted certiorari and affirmed.

\section{Judicial Votes}
Affirmed, 7--2. Stevens wrote for O'Connor, Scalia, Kennedy, Souter, Ginsburg, and Breyer; Rehnquist and Thomas each dissented, joining each other.

\section{Holding}
The right to travel includes the right of those who become permanent residents to be treated like other citizens of their new state, grounded in the Fourteenth Amendment's Citizenship and Privileges or Immunities Clauses. Discrimination by length or place of prior residence gets review ``no less strict'' than strict scrutiny, and Congress cannot authorize it.

\section{Analysis of Opinions}
Stevens noted that even the \textit{Slaughter-House} majority and dissent agreed that a citizen may become a citizen of any state by residence. Deterring poor migrants is impermissible, and fiscal savings, achievable by a 72-cent across-the-board cut, cannot justify 45 subclasses of citizens; the Citizenship Clause forbids ``degrees of citizenship based on length of residence.'' Welfare, unlike portable tuition or divorce, is consumed in-state. Rehnquist replied that a settled resident is ``no longer `traveling' in any sense of the word'' and that the cap was a reasonable residency test. Thomas argued the Clause protects fundamental rights, not welfare benefits, and warned against ``yet another convenient tool for inventing new rights.''

\section{Examples: Future Applications}
\textbf{Same-Side:} capping newcomers' housing stipends at their prior state's level is unconstitutional. \textbf{Opposite-Side:} charging visiting elk hunters ten times the resident fee is valid under Article IV. \textbf{Fence-Sitter:} a two-year wait for in-state tuition applied to adults who have plainly settled.

\section{Critique}
Some originalists praised the Clause's revival, but many echo Thomas that Stevens ignored its history. Laurence Tribe warned that relying on a citizens-only clause could weaken protection for noncitizens, and the revival proved a one-off after \textit{McDonald}.

\section{Key Quotations}
\begin{quote}
``The `right to travel' discussed in our cases embraces at least three different components. It protects the right of a citizen of one State to enter and to leave another State, the right to be treated as a welcome visitor rather than an unfriendly alien when temporarily present in the second State, and, for those travelers who elect to become permanent residents, the right to be treated like other citizens of that State.''
\end{quote}
\noindent --- Justice Stevens, Opinion of the Court, 526 U.S. at 500.

\end{document}
